\begin{center}\textbf{Sommaire}\end{center}
\begin{center}\textbf{Chapitre 1 : Catégories triangulées}\end{center}
\noindent\textbf{\S~1 : Définitions et exemples}\par
\noindent n\textsuperscript{o}~1\quad Définition des catégories triangulées. Premières propriétés.\dotfill 1\par
\noindent n\textsuperscript{o}~2\quad Exemples de catégories triangulées\dotfill 5\par
\noindent n\textsuperscript{o}~3\quad Exemples de foncteurs exacts et de foncteurs cohomologiques.\dotfill 10\par
\medskip\noindent\textbf{\S~2 : Catégories quotients}\par
\noindent n\textsuperscript{o}~1\quad Catégories épaisses. Systèmes multiplicatifs\dotfill 13\par
\noindent n\textsuperscript{o}~2\quad Problèmes universels équivalents\dotfill 15\par
\noindent n\textsuperscript{o}~3\quad Calcul de fractions\dotfill 16\par
\noindent n\textsuperscript{o}~4\quad Propriétés des catégories quotients\dotfill 19\par
\noindent n\textsuperscript{o}~5\quad Propriétés du foncteur de passage au quotient\dotfill 21\par
\noindent n\textsuperscript{o}~6\quad Orthogonalité.\dotfill 24\par
\begin{center}\textbf{Chapitre 2 : Applications aux catégories abéliennes}\end{center}
\noindent\textbf{\S~1 : Les catégories dérivées}\par
\noindent n\textsuperscript{o}~1\quad Les catégories dérivées d'une catégorie abélienne\dotfill 28\par
\noindent n\textsuperscript{o}~2\quad Étude des Ext.\dotfill 34\par
\medskip\noindent\textbf{\S~2 : Les foncteurs dérivés}\par
\noindent n\textsuperscript{o}~1\quad Définitions des foncteurs dérivés\dotfill 37\par
\noindent n\textsuperscript{o}~2\quad Existence des foncteurs dérivés\dotfill 40\par
\noindent n\textsuperscript{o}~3\quad Cup-produit. Composition.\dotfill 44\par
\medskip\noindent\textbf{\S~3 : Exemples}\par
\noindent n\textsuperscript{o}~1\quad Le foncteur $\Hom$\dotfill 46\par
\noindent n\textsuperscript{o}~2\quad $\Tor$ de faisceaux\dotfill 47\par
\noindent n\textsuperscript{o}~3\quad Le foncteur dérivé gauche de l'image réciproque.\dotfill 48\par
